% ======================================================================
\chapter[Тождественные частицы. Принцип Паули. Сложные атомы. Термы]%
        {Тождественные частицы. Принцип Паули. Сложные атомы и атомные термы}
\label{t:13}
% ======================================================================

\begin{qbox}
\textbf{Билет 13.} Тождественность частиц, бозоны и фермионы, принцип Паули. Сложные атомы.
Электронная конфигурация атома. Атомные термы, спектроскопическая запись состояния атома.
\end{qbox}

\section{Содержание билета}

\textbf{Подтемы:}
\begin{enumerate}
 \item принцип неразличимости тождественных частиц; оператор перестановки; симметричные и антисимметричные состояния;
 \item бозоны и фермионы, связь спина со статистикой;
 \item принцип Паули (из антисимметрии), детерминант Слэтера; обменные эффекты;
 \item многоэлектронный атом: приближение самосогласованного поля, оболочки и подоболочки, порядок заполнения;
 \item электронная конфигурация; периодическая система;
 \item $LS$-связь, атомный терм $^{2S+1}L_J$, правила Хунда; примеры.
\end{enumerate}

\textbf{Необходимые понятия:} спин (\tref{t:12}); квантовые числа атома водорода (\tref{t:10});
сложение моментов; квантовый дефект (\tref{t:03}).

\section{Определения и физические принципы}

\begin{principle}{Принцип тождественности (неразличимости)}
Частицы одного сорта (электроны, протоны, фотоны\,\dots) принципиально неразличимы: перестановка двух
тождественных частиц не меняет физического состояния системы. Все наблюдаемые величины симметричны
относительно перестановок.
\end{principle}

\emph{Физическая причина.} В классике одинаковые частицы можно «пометить» по их траекториям. В квантовой
механике траекторий нет: когда волновые пакеты двух электронов перекрываются, невозможно сказать,
какой из них куда попал после взаимодействия.

\begin{definition}[бозоны и фермионы]
\emph{Бозоны} — частицы с целым спином ($0,1,2,\dots$: фотон, $\pi$-мезон, $W^\pm$, $Z^0$, глюон, бозон Хиггса,
ядро $^4$He); волновая функция системы тождественных бозонов \emph{симметрична} относительно перестановки любых
двух частиц. \emph{Фермионы} — частицы с полуцелым спином ($\frac12,\frac32,\dots$: электрон, протон, нейтрон,
нейтрино, кварки, ядро $^3$He); волновая функция \emph{антисимметрична}.
\end{definition}

\begin{principle}{Принцип Паули (1925)}
В системе тождественных фермионов никакие два фермиона не могут находиться в одном и том же
одночастичном квантовом состоянии. Для электронов в атоме: никакие два электрона не могут иметь
одинаковый набор четырёх квантовых чисел $(n,l,m_l,m_s)$.
\end{principle}

\begin{definition}[конфигурация и терм]
\emph{Электронная конфигурация} — распределение электронов по подоболочкам $nl$ (например, $1s^22s^22p^2$).
\emph{Терм} — состояние атома с определёнными полным орбитальным $L$, полным спиновым $S$ (и, после учёта
спин-орбитального взаимодействия, полным $J$) моментами; записывается как $^{2S+1}L_J$, где $L$ обозначается
буквами $S,P,D,F,G,\dots$ для $L=0,1,2,3,4,\dots$; $2S+1$ — \emph{мультиплетность}.
\end{definition}

\emph{Область применимости.} Одночастичное описание (конфигурации) — приближение независимых электронов в
самосогласованном центральном поле. Схема $LS$-связи справедлива для лёгких и средних атомов, где
электростатическое взаимодействие электронов много сильнее спин-орбитального; для тяжёлых атомов
применяется $jj$-связь.

\section{Основная теория}

\subsection{Симметрия волновой функции}

Пусть $\Psi(1,2)$ — волновая функция двух тождественных частиц, где $1\equiv(\vect r_1,\sigma_1)$. Оператор
перестановки $\op P_{12}\Psi(1,2)=\Psi(2,1)$. Из неразличимости $|\Psi(2,1)|^2=|\Psi(1,2)|^2$, т.\,е.
$\Psi(2,1)=\ee^{\ii\alpha}\Psi(1,2)$. Повторная перестановка возвращает исходное состояние:
$\ee^{2\ii\alpha}=1$, $\ee^{\ii\alpha}=\pm1$ (вывод~\ref{der:13:perm}). Итак, возможны лишь два типа состояний:
\begin{equation}\label{eq:13:sym}
  \Psi(2,1)=+\Psi(1,2)\ \ (\text{симметричные}),\qquad \Psi(2,1)=-\Psi(1,2)\ \ (\text{антисимметричные}).
\end{equation}
Гамильтониан тождественных частиц симметричен, поэтому $\op P_{12}$ коммутирует с $\op H$: тип симметрии
сохраняется во времени. Какой тип реализуется, определяет природа частиц.

\begin{principle}{Теорема о связи спина со статистикой (Паули, 1940; без доказательства)}
Частицы с целым спином описываются симметричными волновыми функциями (подчиняются статистике
Бозе--Эйнштейна), с полуцелым — антисимметричными (статистика Ферми--Дирака). Теорема следует из
релятивистской квантовой теории поля; в нерелятивистской квантовой механике она принимается как постулат.
\end{principle}

Составная частица ведёт себя как бозон, если содержит чётное число фермионов, и как фермион — если нечётное
(при условии, что внутренняя структура не возбуждается): $^4$He (2p+2n+2e) — бозон, $^3$He (2p+1n+2e) — фермион.
Это объясняет различие их свойств при низких температурах (сверхтекучесть $^4$He при 2,17~К, у $^3$He — лишь
при $\sim$мК за счёт образования пар).

\subsection{Принцип Паули как следствие антисимметрии}

Для двух невзаимодействующих частиц в одночастичных состояниях $a$ и $b$ (вывод~\ref{der:13:slater})
\begin{equation}\label{eq:13:2p}
  \Psi_{S,A}(1,2)=\frac{1}{\sqrt2}\bigl[\psi_a(1)\psi_b(2)\pm\psi_a(2)\psi_b(1)\bigr].
\end{equation}
При $a=b$ антисимметричная функция тождественно равна нулю — такое состояние невозможно. Это и есть
принцип Паули. Для $N$ фермионов антисимметричная функция — \emph{детерминант Слэтера}
\begin{equation}\label{eq:13:slater}
  \Psi(1,\dots,N)=\frac{1}{\sqrt{N!}}
  \begin{vmatrix}
   \psi_{a_1}(1)&\psi_{a_1}(2)&\cdots&\psi_{a_1}(N)\\
   \vdots&\vdots&&\vdots\\
   \psi_{a_N}(1)&\psi_{a_N}(2)&\cdots&\psi_{a_N}(N)
  \end{vmatrix},
\end{equation}
обращающийся в нуль при совпадении любых двух строк (одинаковые состояния) или столбцов (одинаковые
координаты и спины: два электрона с одинаковым спином не могут оказаться в одной точке).

\paragraph{Обменные эффекты.} Симметрия волновой функции влияет на пространственное распределение даже без
сил между частицами: при $\vect r_1\to\vect r_2$ пространственно-антисимметричная функция обращается в нуль
(частицы «избегают» друг друга), симметричная — усиливается (бозоны «группируются»). Для электронов
полная функция = координатная $\times$ спиновая, и антисимметрична в целом. Если спины параллельны ($S=1$,
симметричная спиновая часть), координатная часть антисимметрична, электроны в среднем дальше друг от друга,
и кулоновское отталкивание меньше. Поэтому энергия зависит от суммарного спина, хотя гамильтониан спина не
содержит, — \emph{обменное взаимодействие}. Оно объясняет первое правило Хунда, расщепление пара- и
ортогелия (у ортогелия, $S=1$, уровни ниже, чем у парагелия с той же конфигурацией) и ферромагнетизм.

\subsection{Сложные атомы: приближение самосогласованного поля}

Гамильтониан атома с $Z$ электронами (ядро неподвижно)
\[
  \op H=\sum_{i=1}^Z\Bigl(-\frac{\hbar^2}{2\mel}\nabla_i^2-\frac{Ze^2}{r_i}\Bigr)+\sum_{i<j}\frac{e^2}{|\vect r_i-\vect r_j|}
\]
не допускает точного решения. В \emph{приближении центрального (самосогласованного) поля} (Хартри--Фок) каждый
электрон движется в усреднённом сферически симметричном поле ядра и остальных электронов $U(r)$: вблизи
ядра $U\approx-Ze^2/r$, вдали $U\approx-e^2/r$ (ядро экранировано $Z-1$ электронами). Одноэлектронные
состояния характеризуются теми же квантовыми числами $n,l,m_l,m_s$, что и в водороде, но поле некулоновское, поэтому
\emph{энергия зависит от $n$ и $l$}: при данном $n$ состояния с меньшим $l$ глубже проникают к ядру, хуже
экранированы и лежат ниже ($E_{ns}<E_{np}<E_{nd}$).

\paragraph{Оболочки и подоболочки.} Совокупность состояний с данным $n$ — \emph{оболочка} ($K,L,M,N,\dots$ для
$n=1,2,3,4,\dots$), с данными $n,l$ — \emph{подоболочка}. По принципу Паули ёмкость подоболочки $2(2l+1)$
($s$ — 2, $p$ — 6, $d$ — 10, $f$ — 14), оболочки — $2n^2$.

\paragraph{Порядок заполнения} (правило Клечковского) — по возрастанию $n+l$, а при равных $n+l$ — по
возрастанию $n$:
\[
  1s\ \ 2s\ \ 2p\ \ 3s\ \ 3p\ \ 4s\ \ 3d\ \ 4p\ \ 5s\ \ 4d\ \ 5p\ \ 6s\ \ 4f\ \ 5d\ \ 6p\ \ 7s\ \ 5f\ \ 6d\ \ 7p .
\]
Например, $4s$ ($n+l=4$) заполняется раньше $3d$ ($n+l=5$): $s$-электрон сильнее проникает к ядру. Правило
эмпирическое и имеет исключения (Cr: $3d^54s^1$, Cu: $3d^{10}4s^1$ — выгодность наполовину или полностью
заполненной $d$-подоболочки).

\paragraph{Электронные конфигурации (примеры).}
\begin{center}
\small
\begin{tabular}{@{}llll@{}}
\toprule
Атом & Конфигурация & Атом & Конфигурация \\
\midrule
H & $1s$ & Na & $1s^22s^22p^63s=[\mathrm{Ne}]3s$ \\
He & $1s^2$ & Cl & $[\mathrm{Ne}]3s^23p^5$ \\
Li & $1s^22s$ & Ar & $[\mathrm{Ne}]3s^23p^6$ \\
C & $1s^22s^22p^2$ & K & $[\mathrm{Ar}]4s$ \\
N & $1s^22s^22p^3$ & Fe & $[\mathrm{Ar}]3d^64s^2$ \\
O & $1s^22s^22p^4$ & Cu & $[\mathrm{Ar}]3d^{10}4s$ \\
Ne & $1s^22s^22p^6$ & Ag & $[\mathrm{Kr}]4d^{10}5s$ \\
\bottomrule
\end{tabular}
\end{center}

\paragraph{Периодичность.} Химические свойства определяются внешними (валентными) электронами. Атомы с
одинаковой структурой внешней подоболочки имеют сходные свойства: щелочные металлы ($ns^1$, легко
отдают электрон, энергия ионизации 4--5~эВ), галогены ($np^5$), инертные газы (заполненные $np^6$, энергия
ионизации максимальна: 21,6~эВ у Ne, 24,6~эВ у He). Длины периодов 2, 8, 8, 18, 18, 32 — следствие порядка
заполнения подоболочек.

\subsection{Атомные термы и $LS$-связь}

Конфигурация не определяет энергию однозначно: остаточное (нецентральное) взаимодействие электронов
расщепляет её на \emph{термы} с разными $L$ и $S$. В схеме \emph{$LS$-связи} (Рассел--Саундерс) складываются отдельно
орбитальные и спиновые моменты:
\begin{equation}\label{eq:13:LS}
  \vect L=\sum_i\vect l_i,\qquad \vect S=\sum_i\vect s_i,\qquad \vect J=\vect L+\vect S,\qquad
  J=|L-S|,\dots,L+S .
\end{equation}
Спин-орбитальное взаимодействие расщепляет терм на компоненты с разными $J$ — \emph{тонкая структура};
при $L\ge S$ число компонент равно мультиплетности $2S+1$ (отсюда название). \textbf{Запись:}
\[
  ^{2S+1}L_J,\qquad\text{например } ^2S_{1/2},\ ^3P_0,\ ^4S_{3/2},\ ^5D_4 .
\]
\textbf{Заполненные подоболочки} имеют $L=0$, $S=0$ (сумма $m_l$ и сумма $m_s$ равны нулю), т.\,е. терм $^1S_0$.
Поэтому терм атома определяется только электронами незаполненных подоболочек.

\begin{principle}{Правила Хунда (для основного терма)}
\begin{enumerate}[leftmargin=1.4em]
 \item Наименьшую энергию имеет терм с максимальным $S$ (при данной конфигурации).
 \item При данном $S$ — терм с максимальным $L$.
 \item Если подоболочка заполнена менее чем наполовину, ниже лежит компонента с $J=|L-S|$; если более чем
       наполовину — с $J=L+S$.
\end{enumerate}
\end{principle}

Первое правило — следствие обменного взаимодействия (параллельные спины $\Rightarrow$ антисимметричная
координатная функция $\Rightarrow$ меньше отталкивание). Второе — электроны, «вращающиеся» в одну сторону,
реже встречаются. Третье — знак константы спин-орбитальной связи для частиц и «дырок».

\paragraph{Основные термы (примеры, см. вывод~\ref{der:13:hund}).}
\begin{center}
\small
\begin{tabular}{@{}llcccl@{}}
\toprule
Атом & Внешняя подоболочка & $S$ & $L$ & $J$ & Терм \\
\midrule
H, Na, Ag & $ns^1$ & $\frac12$ & 0 & $\frac12$ & $^2S_{1/2}$ \\
He, Be & $ns^2$ & 0 & 0 & 0 & $^1S_0$ \\
B & $2p^1$ & $\frac12$ & 1 & $\frac12$ & $^2P_{1/2}$ \\
C & $2p^2$ & 1 & 1 & 0 & $^3P_0$ \\
N & $2p^3$ & $\frac32$ & 0 & $\frac32$ & $^4S_{3/2}$ \\
O & $2p^4$ & 1 & 1 & 2 & $^3P_2$ \\
F & $2p^5$ & $\frac12$ & 1 & $\frac32$ & $^2P_{3/2}$ \\
Fe & $3d^6$ & 2 & 2 & 4 & $^5D_4$ \\
\bottomrule
\end{tabular}
\end{center}

\paragraph{Полная спектроскопическая запись} состояния атома включает конфигурацию и терм, например
основное состояние углерода $1s^22s^22p^2\ ^3P_0$, натрия $[\mathrm{Ne}]3s\ ^2S_{1/2}$, возбуждённое состояние
натрия $[\mathrm{Ne}]3p\ ^2P_{3/2}$. Дублет натрия — переходы $3p\,^2P_{3/2}\to3s\,^2S_{1/2}$ (589,0~нм) и
$3p\,^2P_{1/2}\to3s\,^2S_{1/2}$ (589,6~нм). Правила отбора для дипольных переходов в $LS$-связи: $\Delta S=0$,
$\Delta L=0,\pm1$, $\Delta J=0,\pm1$ (кроме $0\to0$), $\Delta l=\pm1$ для перескакивающего электрона.

\section{Математические выводы}

\begin{derivation}[label=der:13:perm]{собственные значения оператора перестановки}
По определению $\op P_{12}\Psi(1,2)=\Psi(2,1)$. Применив $\op P_{12}$ дважды, получаем исходную функцию:
$\op P_{12}^2=\op 1$. Если $\op P_{12}\Psi=\lambda\Psi$, то $\op P_{12}^2\Psi=\lambda^2\Psi=\Psi$, откуда $\lambda=\pm1$. Состояние
тождественных частиц должно быть собственной функцией $\op P_{12}$ (иначе перестановка давала бы другое,
физически различимое состояние), поэтому оно либо симметрично, либо антисимметрично. Так как
$[\op P_{12},\op H]=0$, тип симметрии сохраняется со временем.
\end{derivation}

\begin{derivation}[label=der:13:slater]{двухчастичные функции и принцип Паули}
Для невзаимодействующих частиц $\op H=\op h(1)+\op h(2)$, и произведение $\psi_a(1)\psi_b(2)$ — решение с энергией
$\varepsilon_a+\varepsilon_b$. Произведение $\psi_a(2)\psi_b(1)$ — тоже решение с той же энергией (обменное вырождение).
Ни одно из произведений не обладает нужной симметрией при $a\neq b$; строим линейные комбинации:
\[
  \Psi_{S,A}=\frac{1}{\sqrt2}\bigl[\psi_a(1)\psi_b(2)\pm\psi_a(2)\psi_b(1)\bigr],\qquad
  \op P_{12}\Psi_{S,A}=\pm\Psi_{S,A}.
\]
Нормировка ($\psi_a\perp\psi_b$): $\int|\Psi|^2=\frac12(1+1\pm2\cdot0)=1$. Антисимметричную функцию можно записать как
определитель $\frac{1}{\sqrt2}\begin{vmatrix}\psi_a(1)&\psi_a(2)\\\psi_b(1)&\psi_b(2)\end{vmatrix}$. При $a=b$ две строки совпадают и $\Psi_A\equiv0$:
\emph{двух фермионов в одном состоянии быть не может}. Для бозонов при $a=b$ $\Psi_S=\psi_a(1)\psi_a(2)$ —
допустимо, и любое число бозонов может находиться в одном состоянии (основа бозе-конденсации, лазеров).
\end{derivation}

\begin{derivation}[label=der:13:hund]{основные термы по правилам Хунда}
\emph{Алгоритм:} размещаем электроны незаполненной подоболочки по ячейкам $m_l=l,l-1,\dots,-l$, сначала по одному
со спином вверх (максимум $S$), начиная с наибольшего $m_l$ (максимум $L$); затем $M_L=\sum m_l=L$,
$M_S=\sum m_s=S$; $J$ — по третьему правилу.

\emph{Углерод $2p^2$:} $m_l=+1\uparrow$, $0\uparrow$. $S=1$, $L=1+0=1$; подоболочка ($6$ мест) заполнена меньше чем наполовину
$\Rightarrow$ $J=|L-S|=0$. Терм $^3P_0$.

\emph{Азот $2p^3$:} $+1\uparrow,0\uparrow,-1\uparrow$. $S=\frac32$, $L=0$, $J=\frac32$. Терм $^4S_{3/2}$.

\emph{Кислород $2p^4$:} $+1\uparrow\downarrow$, $0\uparrow$, $-1\uparrow$. $S=1$, $L=1+1+0-1=1$; больше половины $\Rightarrow$ $J=L+S=2$.
Терм $^3P_2$. (Эквивалентно: две «дырки» в $p$-подоболочке ведут себя как $p^2$ по $L$ и $S$.)

\emph{Железо $3d^6$:} пять электронов $\uparrow$ на $m_l=2,1,0,-1,-2$ (их $\sum m_l=0$) и шестой $\downarrow$ на $m_l=2$.
$S=\frac52-\frac12=2$, $L=2$; больше половины $\Rightarrow$ $J=4$. Терм $^5D_4$.

\emph{Все термы конфигурации $np^2$} (без доказательства полного перебора): из $\binom62=15$ состояний, допускаемых
принципом Паули, образуются термы $^1S$ (1 состояние), $^1D$ (5) и $^3P$ (9): $1+5+9=15$. Терма $^3D$ нет — он потребовал
бы двух электронов с $m_l=1$ и одинаковыми спинами. Для неэквивалентных электронов ($2p3p$) возможны все
комбинации $L=0,1,2$, $S=0,1$.
\end{derivation}

\begin{derivation}[label=der:13:capacity]{ёмкость подоболочки и оболочки}
При данных $n,l$ квантовое число $m_l$ принимает $2l+1$ значений, $m_s$ — два. По принципу Паули каждый набор
$(n,l,m_l,m_s)$ занят не более чем одним электроном, поэтому ёмкость подоболочки $2(2l+1)$, оболочки
$\sum_{l=0}^{n-1}2(2l+1)=2n^2$ (2, 8, 18, 32).
\end{derivation}

\section{Итоговые формулы}

\keyformula{\Psi(\dots,i,\dots,j,\dots)=\pm\Psi(\dots,j,\dots,i,\dots)}
 {$+$ — бозоны (целый спин), $-$ — фермионы (полуцелый спин); $i\equiv(\vect r_i,\sigma_i)$.}
 {любая система тождественных частиц.}
 {симметрия разрешённых состояний; принцип Паули; статистика.}

\keyformula{\Psi_A=\frac{1}{\sqrt{N!}}\det\bigl[\psi_{a_k}(i)\bigr],\qquad N_{nl}^{\max}=2(2l+1),\quad N_n^{\max}=2n^2}
 {$\psi_{a_k}$ — одночастичные спин-орбитали.}
 {невзаимодействующие фермионы (приближение самосогласованного поля).}
 {принцип Паули; ёмкость подоболочек и оболочек.}

\keyformula{^{2S+1}L_J,\qquad \vect L=\sum\vect l_i,\ \ \vect S=\sum\vect s_i,\ \ J=|L-S|,\dots,L+S}
 {$L$ — полный орбитальный момент ($S,P,D,F$ для $0,1,2,3$), $S$ — полный спин, $J$ — полный момент.}
 {$LS$-связь (лёгкие и средние атомы).}
 {классификация состояний атома; основной терм по правилам Хунда.}

\section{Физическая интерпретация}

\textbf{Без принципа Паули} все электроны оказались бы в состоянии $1s$, атомы разных элементов различались бы
лишь размерами, не было бы периодической системы и химии; вещество было бы неустойчиво к сжатию.
Принцип Паули объясняет оболочечную структуру атомов, устойчивость вещества, свойства металлов
(электронный газ Ферми), белых карликов и нейтронных звёзд (давление вырожденного газа).

\textbf{Энергетическая роль спина.} Гамильтониан без спин-орбитальных членов спина не содержит, но
допустимая симметрия координатной части зависит от суммарного спина. Поэтому энергия термов с разными $S$
различается на величину порядка электростатической (эВ) — много больше магнитной энергии взаимодействия спинов.

\textbf{Эксперимент.} Схема термов проверяется по спектрам: мультиплетность линий, правила отбора, расщепление
в магнитном поле с фактором Ланде $g_J$ (\tref{t:12}). Рентгеновские спектры (закон Мозли) подтверждают оболочечную
структуру внутренних электронов.

\section{Что необходимо знать на экзамене}

\begin{examcheck}
\checkhead{Определения}
\begin{itemize}
 \item принцип тождественности; бозоны, фермионы; связь спина со статистикой; принцип Паули;
 \item приближение самосогласованного поля; оболочка, подоболочка, конфигурация; терм, мультиплетность, $LS$-связь.
\end{itemize}
\checkhead{Формулы}
\begin{itemize}
 \item $\Psi_{S,A}=\frac{1}{\sqrt2}[\psi_a(1)\psi_b(2)\pm\psi_a(2)\psi_b(1)]$; детерминант Слэтера;
 \item ёмкости $2(2l+1)$ и $2n^2$; правило $n+l$; $^{2S+1}L_J$, $J=|L-S|..L+S$.
\end{itemize}
\checkhead{Выводы}
\begin{itemize}
 \item собственные значения $\op P_{12}$ равны $\pm1$; принцип Паули из антисимметрии;
 \item основные термы C, N, O, Fe по правилам Хунда; термы $np^2$.
\end{itemize}
\checkhead{Интерпретации}
\begin{itemize}
 \item обменное взаимодействие и первое правило Хунда; орто- и парагелий;
 \item зависимость $E$ от $l$ в сложных атомах (проникновение, экранирование); периодическая система; исключения Cr, Cu.
\end{itemize}
\end{examcheck}
